% !TEX root = ../main.tex
% !TEX program = xelatex
\section{模型评价与推广}

\subsection{模型优点与创新}

\begin{enumerate}
    \item \textbf{以实际费用检验预测价值。}
    预测与调度共同服务于购电节费。近期加权准备较等权方案净节省44.66万元；日内价格预测则兼顾费用与紧急补购，使方法选择落实到运行代价。

    \item \textbf{单独量化新增信息的节费贡献。}
    同一模型下，日内反馈与改单节省50.63万元，后发光伏预报再节省6.95万元。对照区分了两类收益，且新预报的节费并非来自多消耗期末电量。

    \item \textbf{统一储能约束，支持连续运行与核验。}
    四问共用电量平衡与储能约束，逐10分钟核对库存衔接和设备限制，并复算账单，便于核验连续运行与节费结果。
\end{enumerate}

\subsection{模型局限与改进方向}

现有参数比较基于2025年数据，后续可通过跨年验证完善校准，补充供需与电价共同异常的情景，并纳入电池老化和网络损耗成本，使方案更贴近长期运行条件。

\subsection{模型推广性评估}

该框架可为校园、社区和工业园区微网提供参考。按当地供需、设备和交易规则校准后，可扩展至风光互补、多台储能及可调负载，协调本地发电与外网购电，兼顾用能经济性和可再生能源利用。
